% Options for packages loaded elsewhere
\PassOptionsToPackage{unicode}{hyperref}
\PassOptionsToPackage{hyphens}{url}
\documentclass[
]{article}
\usepackage{xcolor}
\usepackage{amsmath,amssymb}
\setcounter{secnumdepth}{-\maxdimen} % remove section numbering
\usepackage{iftex}
\ifPDFTeX
  \usepackage[T1]{fontenc}
  \usepackage[utf8]{inputenc}
  \usepackage{textcomp} % provide euro and other symbols
\else % if luatex or xetex
  \usepackage{unicode-math} % this also loads fontspec
  \defaultfontfeatures{Scale=MatchLowercase}
  \defaultfontfeatures[\rmfamily]{Ligatures=TeX,Scale=1}
\fi
\usepackage{lmodern}
\ifPDFTeX\else
  % xetex/luatex font selection
\fi
% Use upquote if available, for straight quotes in verbatim environments
\IfFileExists{upquote.sty}{\usepackage{upquote}}{}
\IfFileExists{microtype.sty}{% use microtype if available
  \usepackage[]{microtype}
  \UseMicrotypeSet[protrusion]{basicmath} % disable protrusion for tt fonts
}{}
\makeatletter
\@ifundefined{KOMAClassName}{% if non-KOMA class
  \IfFileExists{parskip.sty}{%
    \usepackage{parskip}
  }{% else
    \setlength{\parindent}{0pt}
    \setlength{\parskip}{6pt plus 2pt minus 1pt}}
}{% if KOMA class
  \KOMAoptions{parskip=half}}
\makeatother
\setlength{\emergencystretch}{3em} % prevent overfull lines
\providecommand{\tightlist}{%
  \setlength{\itemsep}{0pt}\setlength{\parskip}{0pt}}
\ifPDFTeX
  \TeXXeTstate=1
  \newcommand{\RL}[1]{\beginR #1\endR}
  \newcommand{\LR}[1]{\beginL #1\endL}
  \newenvironment{RTL}{\beginR}{\endR}
  \newenvironment{LTR}{\beginL}{\endL}
\fi
\usepackage{bookmark}
\IfFileExists{xurl.sty}{\usepackage{xurl}}{} % add URL line breaks if available
\urlstyle{same}
\hypersetup{
  pdftitle={2. Невязкое несжимаемое течение},
  hidelinks,
  pdfcreator={LaTeX via pandoc}}

\title{2. Невязкое несжимаемое течение}
\author{}
\date{}

\begin{document}
\maketitle

\subsection{Тензор неоднородности скорости среды и его
декомпозиция}\label{ux442ux435ux43dux437ux43eux440-ux43dux435ux43eux434ux43dux43eux440ux43eux434ux43dux43eux441ux442ux438-ux441ux43aux43eux440ux43eux441ux442ux438-ux441ux440ux435ux434ux44b-ux438-ux435ux433ux43e-ux434ux435ux43aux43eux43cux43fux43eux437ux438ux446ux438ux44f}

Рассмотрим движение произвольного элемента среды. Для простоты будем
рассматривать плоскую задачу.\\
Полный тензор, характеризующий неоднородность поля скорости, запишется
как

Как известно, любой тензор можно представить в виде суммы симметричной и
антисимметричной части\\
{} Тензор S носит название тензора скоростей деформации (S). Тензор {}
отвечает за локальное вращение среды как "твердого тела".

Компоненты этого тензора содержат коvпоненты вектора завихренности -
ротора скорости среды

\subsection{Уравнение Фридмана (уравнение генерации и переноса
завихренности)}\label{ux443ux440ux430ux432ux43dux435ux43dux438ux435-ux444ux440ux438ux434ux43cux430ux43dux430-ux443ux440ux430ux432ux43dux435ux43dux438ux435-ux433ux435ux43dux435ux440ux430ux446ux438ux438-ux438-ux43fux435ux440ux435ux43dux43eux441ux430-ux437ux430ux432ux438ux445ux440ux435ux43dux43dux43eux441ux442ux438}

Если взять от уравнения Эйлера ротор, получим\\
{}

Подставляя сюда выражение для завихренности и раскрывая скобки,
получим\\
{}

Первые два члена в этом выражении представляют собой лагранжеву
производную от завихренности. Таким образом, изменение завихренности в
некотором объеме обусловлено деформацией и растяжением вихря (два члена
справа), а также внешними факторами: бароклинностью (неколлинеарностью
градиентов давления и плотности), а также непотенциальностью массовой
силы, действующей на среду.

Отдельный интерес представляет ситуация, когда внешняя сила потенциальна
(например, в гравитационном поле), а жидкость покоится. В этом случае
можно видеть, что равновесие среды возможно только при {}, или когда
давление и плотность однозначно связаны {}. Такая ситуация реализуется в
случае, когда жидкость несжимаемая (плотность постоянна), или же когда
термодинамический процесс в среде задан однозначно и представляет собой,
например, изотерму или адиабату. В таких случаях жидкость называется
\textbf{\emph{баротропной}}. Если это соотношение не выполняется
(например, есть градиент температуры), среда называется
\textbf{\emph{бароклинной}}.

Таким образом, в \textbf{\emph{баротропной}} жидкости в поле
потенциальных сил внешние источники завихренности отсутствуют.

Задача : гидростатическое равновесие баротропной жидкости

\textbf{1. Пусть жидкость покоится в поле сил тяжести, температура среды
постоянна. Найти распределение плотности.}\\
Уравнение гидростатического равновесия получим из уравнения Эйлера\\
{}При постоянной температуре {}.\\
Тогда имеем {}{}

\subsection{Безвихревое (потенциальное) движение
среды}\label{ux431ux435ux437ux432ux438ux445ux440ux435ux432ux43eux435-ux43fux43eux442ux435ux43dux446ux438ux430ux43bux44cux43dux43eux435-ux434ux432ux438ux436ux435ux43dux438ux435-ux441ux440ux435ux434ux44b}

Покажем, что если в начальный момент движение идеальной жидкости
является безвихревым, то оно останется таковым и впредь. Для этого
покажем, что циркуляция скорости по любому замкнутому контуру.
вмороженному в среду, сохраняется.\\
Проинтегрируем поле скорости в нашем объеме,

{}В этом выражении последний член равен нулю (можно проверить
интегрированием по частям). В то же время подынтегральное выражение в
первом члене в правой части может быть записан из уравнения Эйлера как\\
{}Если сила в правой части уравнения потенциальна, то интеграл от правой
части равен нулю (градиент любой скалярной функции- потенциальная
функция).\\
Тогда имеем\\
{}Таким образом, \textbf{\emph{в идеальной баротропной жидкости в поле
потенциальных сил циркуляция скорости по любому замкнутому контуру
сохраняется.}}\\
В этом выражении контурный интеграл можно заменить на интеграл по
поверхности\\
{} и потребовать его выполнения для бесконечно малого контура, получим,
что \textbf{\emph{если движение идеальной баротропной жидкости во всем
объеме было в некоторый момент времени безвихревым, оно останется таким
в любой последующий момент времени}}.

\subsection{Потенциал
скорости}\label{ux43fux43eux442ux435ux43dux446ux438ux430ux43b-ux441ux43aux43eux440ux43eux441ux442ux438}

Для невязкого несжимаемого течения уравнения неразрывности запишется как
{}

В тех областях, где течение безвихревое, мы можем ввести скалярную
величину- потенциал скорости {} ,тогда\\
{}Таким образом, уравнение неразрывности для несжимаемой может быть, с
повышением порядка уравнения, сведено к уравнению Лапласа для
потенциала.\\
В силу этого факта безвихревое течение называют иногда еще
\textbf{\emph{потенциальным}} течением.\\
Важно, что потенциал скорости однозначно определен только для случая,
когда течение безвихревое во всем исследуемом объеме. Если это правило
нарушается по меньшей мере в некоторых точках пространства, то потенциал
становится многозначным и его определение зависит от конкретного контура
интегрирования (количества обходов областей с ненулевой завихренностью).

\subsection{Потенциал скорости и уравнение Бернулли в потенциальном
течении}\label{ux43fux43eux442ux435ux43dux446ux438ux430ux43b-ux441ux43aux43eux440ux43eux441ux442ux438-ux438-ux443ux440ux430ux432ux43dux435ux43dux438ux435-ux431ux435ux440ux43dux443ux43bux43bux438-ux432-ux43fux43eux442ux435ux43dux446ux438ux430ux43bux44cux43dux43eux43c-ux442ux435ux447ux435ux43dux438ux438}

Подставляя определение потенциала в уравнение Эйлера, и используя
соотношение\\
{}получим в безвихревой области\\
{}Или, интегрируя это выражение по пространству, {}Это выражение можно
рассматривать как нестационарный аналог уравнения Бернулли

\section{Двумерное плоское
течение}\label{ux434ux432ux443ux43cux435ux440ux43dux43eux435-ux43fux43bux43eux441ux43aux43eux435-ux442ux435ux447ux435ux43dux438ux435}

\subsection{Функция
тока}\label{ux444ux443ux43dux43aux446ux438ux44f-ux442ux43eux43aux430}

Для плоских и осесимметричных течений полезным является введение функции
тока. Введем скалярную функцию {}Тогда уравнение неразрывности
удовлетворится автоматически, ведь\\
{}Для стационарного течения постоянное значение функции тока сохраняется
вдоль любой из линий тока. Действительно, уравнение линии тока можем
записать как\\
{}подставляя выражение для функции тока, получим

{}Таким образом, \textbf{\emph{в стационарном двумерном течении
несжимаемой жидкости функция тока остается постоянной вдоль линии
тока}}\\
Изоповерхности потенциала скорости и линии тока перпендикулярны друг
другу в каждой точке пространства, в этом смысле есть аналогия с
потенциалом и напряженностью электрического поля в электростатике.

\textbf{Задача}

Использование уравнения Лапласа для вычисления поля скорости у шара в
однородном потоке несжимаемой жидкости\\
Постановка задачи показана на рис. На цилиндр слева направо натекает
поток со скоростью {}\\
На бесконечности скорость равна скорости невозмущенного течения, на
поверхности шара {}\\
Наша цель- найти поле скоростей в пространстве, окружающем шар,
посчитать давление на его поверхности и силу сопротивления\\
\textbf{Решение}\\
Потенциал свободного однородного течения очевидно может быть найден как
{}.\\
Поскольку задача линейная, мы можем представить течение вокруг шара в
виде суммы свободного потока и возмущения: {}\\
Сформулируем задачу для возмущения потенциала, связанного с присутсnвием
тела:

Собственными функциями уравнения Лапласа в трехмерной задаче, убывающими
на бесконечности, являются {}. Поскольку наше решение должно линейно
зависеть от {}, мы будем искать его в виде {}На поверхности цилиндра
нормальная компонента скорости равна 0, поэтому {}{}Поле скорости
запишется как

Линии тока и потенциал для этого течения построены на рис.\\
На поверхности шара скорость потока определяется касательной компонентой
и равна {}График скорости показан на рис. Видно, что есть две особых
точки расположенных здесь симметрично, в которых скорость обращается в
ноль- точка присоединения потока и точка схода потока в кормовой области
тела.\\
Из уравнения Бернулли найдем давление:\\
{}Давление максимально в лобовой и кормовой точке, и минимально в
миделевом сечении шара.\\
Сила сопротивления шара равна {}. Поскольку распределение давления
симметрично относительно миделева сечения, {}. Этот вывод является
частным случаем так называемого парадокса Д\textquotesingle Аламбера -
отсутствия сопротивления тел при обтекании их идеальной несжимаемой
средой.

\textbf{\emph{Парадокс Д\textquotesingle Алабмера - отсутствие
сопротивления у тел при обтекании их идеальной несжимаемой средой}}
можно объяснить так:\\
Окружим наше тело контрольной поверхностью.\\
Проинтегрируем внутри нее стационарное уравнение Эйлера, предварительно
умножив его на скорость набегающего потока {}.\\
{}В правой части- работа силы сопротивления,\\
{}Левую перепишем как {}В безвихревой области\\
{} {}Вдоль любой линии тока величина {} постоянна, так как является
интегралом Бернулли. Для однородного набегающего потока можем вообще
вынести ее за скобки, после чего левая часть оказывается равна нулю. В
общем же случае данное выражение представляет собой поток полной
кинетической энергии среды. Если разбить объем на струйки тока (рис),
вдоль каждой из них поток должен сохраняться, так как течения идеальной
среды не приводит к диссипации кинетической энергии. Тогда поток ее по
замкнутой поверхности должен быть равен нулю, а значит равно нулю
сопротивление тела обтеканию.

\end{document}
